\documentclass[12pt,a4paper]{article}

% ---------- Paquetes ----------
\usepackage[utf8]{inputenc}
\usepackage[T1]{fontenc}
\usepackage[spanish]{babel}
\usepackage[margin=2.5cm, top=3cm, bottom=3cm]{geometry}
\usepackage{hyperref}
\usepackage{graphicx}
\usepackage{booktabs}
\usepackage{xcolor}
\usepackage{listings}
\usepackage{tikz}
\usepackage{float}
\usepackage{parskip}
\usepackage{fancyhdr}
\usepackage{array}
\usepackage{mdframed}
\usepackage{lmodern}
\usepackage[scaled=0.95]{helvet}
\usepackage{courier}
\renewcommand{\familydefault}{\sfdefault}
\usetikzlibrary{shapes.geometric, arrows.meta, positioning, fit, backgrounds}

% ---------- Colores de marca nexo ----------
\definecolor{nexoBlack}{HTML}{0A0E1A}
\definecolor{nexoCyan}{HTML}{00D4FF}
\definecolor{nexoMidnight}{HTML}{1A2744}
\definecolor{nexoGhost}{HTML}{F4F7FC}
\definecolor{nexoSlate}{HTML}{8A95A8}
\definecolor{nexoGreen}{HTML}{00E5A0}
\definecolor{nexoAmber}{HTML}{FFB800}
\definecolor{nexoRed}{HTML}{FF4D6A}
\definecolor{codeBg}{HTML}{F4F7FC}

% ---------- Listings ----------
\lstset{
  backgroundcolor=\color{codeBg},
  basicstyle=\ttfamily\footnotesize,
  keywordstyle=\color{nexoCyan!80!black}\bfseries,
  stringstyle=\color{nexoAmber!80!black},
  commentstyle=\color{nexoSlate}\itshape,
  breaklines=true,
  frame=single,
  framesep=5pt,
  rulecolor=\color{nexoSlate!30},
  showstringspaces=false,
  numbers=left,
  numberstyle=\tiny\color{nexoSlate},
  numbersep=8pt,
  tabsize=4,
  xleftmargin=1.5em,
}

% ---------- Encabezado / pie de página ----------
\pagestyle{fancy}
\fancyhf{}
\renewcommand{\headrulewidth}{0pt}
\fancyhead[L]{%
  \begin{tikzpicture}[remember picture, overlay]
    \fill[nexoBlack] (current page.north west)
      rectangle ([yshift=-1.1cm]current page.north east);
    \node[anchor=west, text=white, font=\bfseries\small]
      at ([xshift=2.5cm, yshift=-0.55cm]current page.north west)
      {\textcolor{nexoCyan}{nexo} \textcolor{nexoSlate}{---}\ Infraestructura Frontend};
    \node[anchor=east, text=nexoSlate, font=\small]
      at ([xshift=-2.5cm, yshift=-0.55cm]current page.north east)
      {Documento Interno · v2.0};
  \end{tikzpicture}%
}
\fancyfoot[C]{\small\textcolor{nexoSlate}{\thepage}}

% ---------- Hyperref ----------
\hypersetup{
  colorlinks=true,
  linkcolor=nexoCyan!60!black,
  urlcolor=nexoCyan!60!black,
  pdftitle={Infraestructura Frontend — nexo courier},
  pdfauthor={nexo Engineering},
}

% ---------- TikZ ----------
\tikzset{
  block/.style={
    rectangle, rounded corners=4pt, minimum width=2.8cm, minimum height=1cm,
    text centered, draw=nexoMidnight, fill=nexoMidnight!15, font=\small
  },
  aws/.style={
    rectangle, rounded corners=4pt, minimum width=2.8cm, minimum height=1cm,
    text centered, draw=nexoCyan!70!black, fill=nexoCyan!10, font=\small\bfseries
  },
  external/.style={
    ellipse, minimum width=2.5cm, minimum height=0.9cm,
    text centered, draw=nexoSlate, fill=nexoSlate!10, font=\small
  },
  arrow/.style={-{Stealth[length=6pt]}, thick, color=nexoMidnight},
  cyan_arrow/.style={-{Stealth[length=6pt]}, thick, color=nexoCyan!70!black},
  dashed_arrow/.style={-{Stealth[length=6pt]}, dashed, thick, color=nexoSlate},
}

% ---------- Caja nexo ----------
\newmdenv[
  backgroundcolor=nexoBlack!5,
  linecolor=nexoCyan,
  linewidth=2pt,
  leftline=true, rightline=false, topline=false, bottomline=false,
  innerleftmargin=12pt, innerrightmargin=8pt,
  innertopmargin=6pt, innerbottommargin=6pt,
]{nexobox}

% ==========================================================
\begin{document}

% ---------- PORTADA ----------
\begin{titlepage}
  \pagecolor{nexoBlack}
  \color{white}
  \centering
  \vspace*{3cm}

  {\large\textcolor{nexoSlate}\itshape Tu puente entre dos mundos.\par}
  \vspace{1.2cm}

  {\fontsize{52}{56}\selectfont\bfseries
    \textcolor{white}{ne}\textcolor{nexoCyan}{x}\textcolor{white}{o}\par}
  \vspace{0.4cm}
  {\LARGE\bfseries Frontend --- Arquitectura Técnica\par}

  \vspace{1.8cm}
  \textcolor{nexoCyan}{\rule{0.55\textwidth}{1.5pt}}%
  \textcolor{nexoCyan}{\large\ $\bullet$}\par

  \vspace{1.5cm}
  {\large
    \textcolor{nexoSlate}{Empresa}\quad\textcolor{white}{nexo courier (Costa Rica)} \\[0.5em]
    \textcolor{nexoSlate}{Sistema}\quad\textcolor{white}{Sitio web y panel de administración} \\[0.5em]
    \textcolor{nexoSlate}{Framework}\quad\textcolor{white}{Next.js 14 (App Router)} \\[0.5em]
    \textcolor{nexoSlate}{Hosting}\quad\textcolor{white}{AWS Amplify Hosting} \\[0.5em]
    \textcolor{nexoSlate}{Dominio}\quad\textcolor{white}{nexocourier.com} \\[0.5em]
    \textcolor{nexoSlate}{Versión}\quad\textcolor{white}{2.0} \\[0.5em]
    \textcolor{nexoSlate}{Fecha}\quad\textcolor{white}{Abril 2026}
  }

  \vfill
  {\small\textcolor{nexoSlate}{Confidencial --- Uso Interno}}
\end{titlepage}

\pagecolor{white}
\color{black}

\tableofcontents
\newpage

% ==========================================================
\section{Resumen Ejecutivo}

\begin{nexobox}
El frontend de \textbf{nexo courier} es una aplicación web moderna construida con
\textbf{Next.js 14}. Se despliega automáticamente en \textbf{AWS Amplify Hosting}
cada vez que se hace push a la rama \texttt{master} de GitHub. El sitio vive en
\textbf{nexocourier.com} con HTTPS gestionado por AWS Certificate Manager.
\end{nexobox}

\vspace{0.5em}
El sistema incluye un sitio público de marketing, un portal de casillero para clientes
autenticados, un sistema de gestión de pedidos con seguimiento de estado, y un panel de
administración protegido por grupo de Cognito. La autenticación se maneja directamente
contra \textbf{AWS Cognito} usando el SDK de Amplify v6, con verificación obligatoria
de correo electrónico al registrarse.

\subsection{Stack tecnológico}

\begin{table}[H]
\centering
\begin{tabular}{@{}lll@{}}
\toprule
\textbf{Componente} & \textbf{Tecnología} & \textbf{Versión} \\
\midrule
Framework & Next.js (App Router) & 14.2.35 \\
Lenguaje & TypeScript & \textasciicircum{}5 \\
Estilos & Tailwind CSS & \textasciicircum{}3.4.17 \\
Auth SDK & AWS Amplify & \textasciicircum{}6.16.4 \\
Formularios & React Hook Form + Zod & \textasciicircum{}7.54 / \textasciicircum{}3.23 \\
Animaciones & Framer Motion & \textasciicircum{}11.11 \\
Iconos & Lucide React & \textasciicircum{}0.468 \\
Hosting & AWS Amplify Hosting & -- \\
DNS & Cloudflare & -- \\
\bottomrule
\end{tabular}
\caption{Stack tecnológico del frontend}
\end{table}

% ==========================================================
\section{Pipeline CI/CD}

\begin{lstlisting}[caption={Pipeline de despliegue automático}, numbers=none]
  Developer        GitHub                  AWS Amplify        CloudFront
  (local)          nexo-git/nexocode        Build              CDN
  +---------+      +------------------+    +------------+    +----------+
  |  code   | ---> |  master branch   | -> |  npm build | -> |  deploy  |
  +---------+      +------------------+    +------------+    +----------+
  git push          webhook trigger         ~3 minutos         nexocourier.com
\end{lstlisting}

\subsection*{Fases del build}

\begin{description}
  \item[\textbf{preBuild}] \texttt{npm ci --legacy-peer-deps} --- instala dependencias exactas del lock file.
  \item[\textbf{build}] \texttt{npm run build} --- compila Next.js y genera artefactos en \texttt{.next/}.
  \item[\textbf{deploy}] Amplify toma el directorio \texttt{.next/} y lo publica en CloudFront globalmente.
\end{description}

% ==========================================================
\section{Estructura de Archivos}

\begin{lstlisting}[language=bash, caption={Estructura del proyecto}]
nexocode/
├── src/
│   ├── app/                       # App Router (Next.js 14)
│   │   ├── layout.tsx             # Layout raiz + AmplifyProvider
│   │   ├── page.tsx               # Landing page publica
│   │   ├── casillero/page.tsx     # Casillero USA (registro/info)
│   │   ├── login/page.tsx         # Inicio de sesion
│   │   ├── cuenta/page.tsx        # Perfil del usuario autenticado
│   │   ├── pedidos/page.tsx       # Historial de pedidos (protegida)
│   │   ├── cotizar/page.tsx       # Calculadora de envios
│   │   ├── asistente/page.tsx     # Asistente de compras IA
│   │   ├── nosotros/page.tsx      # Pagina institucional
│   │   ├── tracking/page.tsx      # Rastreo de paquetes
│   │   └── admin/
│   │       ├── page.tsx           # Panel admin: usuarios + pedidos
│   │       └── [id]/page.tsx      # Editar/eliminar usuario
│   ├── components/
│   │   ├── AmplifyProvider.tsx    # Init del SDK de Amplify (nivel modulo)
│   │   ├── ui/                    # Button, Card, Select, Badge, Input
│   │   ├── casillero/             # CasilleroForm, LoginForm, AddressCard
│   │   ├── layout/                # Navbar (con dropdown), Footer
│   │   ├── cotizar/               # QuoteForm, QuoteResult
│   │   ├── tracking/              # StatusBadge, TrackingCard, TrackingTimeline
│   │   └── sections/              # Secciones de la landing
│   ├── lib/
│   │   ├── casillero.ts           # Auth: registro, login, verificacion, perfil
│   │   ├── orders.ts              # Pedidos: crear, listar, actualizar
│   │   ├── amplify-config.ts      # Configuracion Amplify SDK
│   │   ├── constants.ts           # Constantes globales
│   │   ├── pricing.ts             # Logica de cotizacion
│   │   ├── tracking.ts            # Integracion de rastreo
│   │   └── utils.ts               # Utilidades (cn, etc.)
│   └── types/
│       ├── casillero.ts           # NexoUser, NexoOrder, OrderStatus
│       ├── quote.ts               # Tipos de cotizacion
│       └── tracking.ts            # Tipos de rastreo
├── public/                        # Assets estaticos
├── docs/                          # Documentacion tecnica (este archivo)
├── amplify.yml                    # Config build Amplify
├── next.config.mjs
├── tailwind.config.ts
├── tsconfig.json
└── package.json
\end{lstlisting}

% ==========================================================
\section{Páginas y Rutas}

\begin{table}[H]
\centering
\begin{tabular}{@{}lll@{}}
\toprule
\textbf{Ruta} & \textbf{Acceso} & \textbf{Descripción} \\
\midrule
\texttt{/} & Público & Landing page con secciones de marketing \\
\texttt{/cotizar} & Público & Calculadora de costos de envío \\
\texttt{/asistente} & Público & Asistente de compras con IA \\
\texttt{/nosotros} & Público & Página institucional \\
\texttt{/tracking} & Público & Rastreo de paquetes por tracking number \\
\texttt{/casillero} & Público / Auth & Dirección USA si hay sesión activa \\
\texttt{/login} & Público & Formulario de inicio de sesión \\
\texttt{/pedidos} & Protegida & Historial y gestión de pedidos del usuario \\
\texttt{/cuenta} & Protegida & Editar perfil y eliminar cuenta \\
\texttt{/admin} & Protegida (admin) & Dashboard: tabs Usuarios y Pedidos \\
\texttt{/admin/[id]} & Protegida (admin) & Editar o eliminar un usuario \\
\bottomrule
\end{tabular}
\caption{Rutas del sitio web}
\end{table}

\begin{nexobox}
\textbf{Protección de rutas:} No existe middleware de Next.js. Cada página protegida
verifica la sesión en su propio \texttt{useEffect} usando \texttt{getCurrentUser()}.
Las páginas de admin además verifican que el usuario pertenezca al grupo \texttt{admin}
de Cognito leyendo el claim \texttt{cognito:groups} del JWT.
\end{nexobox}

% ==========================================================
\section{Autenticación con AWS Amplify SDK}

\subsection{Flujo de registro con verificación de correo}

\begin{lstlisting}[caption={Flujo de registro con verificación de correo (double opt-in)}, numbers=none]
  1.  Usuario completa el formulario de registro
       |
  2.  registerUser() llama a signUp() de Amplify
       |
  3.  Cognito crea la cuenta y envia codigo de 6 digitos al correo
       |
  4.  registerUser() retorna { needsConfirmation: true, email }
       |
  5.  CasilleroForm muestra pantalla de ingreso del codigo
       |
  6.  Usuario ingresa el codigo de verificacion
       |
  7.  confirmAndLogin(email, password, code)
        -> confirmSignUp()  [valida el codigo en Cognito]
        -> signIn()         [inicia sesion automaticamente]
       |
  8.  Frontend redirige a /casillero
\end{lstlisting}

\subsection{Flujo de login}

\begin{lstlisting}[caption={Flujo de inicio de sesión}, numbers=none]
  1.  Usuario ingresa email y contrasena
       |
  2.  loginUser() llama a signIn() de Amplify
       |
  3.  Cognito valida credenciales y retorna tokens JWT
       |
  4.  SDK almacena tokens en localStorage del navegador
       |
  5.  getCurrentUser() -> fetchUserAttributes() -> NexoUser
       |
  6.  Frontend redirige a /casillero
\end{lstlisting}

\begin{nexobox}
\textbf{Amplify v6 usa \texttt{localStorage}, no cookies.} Por eso no existe middleware
de Next.js en este proyecto: el middleware corre en Edge Runtime y no tiene acceso a
\texttt{localStorage}. La protección de rutas se hace client-side en cada página.
\end{nexobox}

\subsection{Funciones de auth --- src/lib/casillero.ts}

\begin{table}[H]
\centering
\begin{tabular}{@{}ll@{}}
\toprule
\textbf{Función} & \textbf{Descripción} \\
\midrule
\texttt{registerUser(data)} & Crea cuenta en Cognito. Retorna \texttt{\{needsConfirmation, email\}} \\
\texttt{confirmAndLogin(email, password, code)} & Confirma código y hace login automático \\
\texttt{resendCode(email)} & Reenvía el código de verificación al correo \\
\texttt{loginUser(email, password)} & Inicia sesión y retorna \texttt{NexoUser} \\
\texttt{getCurrentUser()} & Retorna el usuario autenticado o \texttt{null} \\
\texttt{updateCurrentUser(data)} & Actualiza nombre, apellido, móvil, teléfono \\
\texttt{deleteCurrentUser()} & Elimina la cuenta de Cognito \\
\texttt{logoutUser()} & Cierra la sesión activa \\
\texttt{generateAddress(nombre, apellido)} & Genera dirección de casillero USA \\
\bottomrule
\end{tabular}
\caption{API de autenticación --- src/lib/casillero.ts}
\end{table}

\subsection{Tipos principales --- src/types/casillero.ts}

\begin{lstlisting}[language=bash, caption={Tipos del dominio principal}]
// Usuario autenticado
interface NexoUser {
  id:        string            // sub de Cognito
  tipo:      "persona" | "empresa"
  nombre:    string
  apellido:  string
  email:     string
  movil:     string
  telefono:  string
  createdAt: string
}

// Pedido de envio
type OrderStatus = "en_ruta" | "atascado" | "entregado"

interface NexoOrder {
  orderId:       string        // UUID
  userId:        string        // sub de Cognito del dueno
  userName:      string
  userEmail:     string
  trackingNumber: string
  description:   string
  startDate:     string        // ISO 8601
  status:        OrderStatus
  peso?:         number        // libras, lo carga el admin
  totalPagado?:  number        // USD, lo carga el admin
  updatedAt:     string
}
\end{lstlisting}

% ==========================================================
\section{Sistema de Pedidos}

\subsection{Descripción general}

Los clientes pueden registrar sus pedidos con número de tracking y descripción.
El admin ve todos los pedidos, actualiza el estado y carga el peso y total cobrado.

\begin{lstlisting}[caption={Flujo del sistema de pedidos}, numbers=none]
  Cliente (/pedidos)
    GET /orders   -->  API Gateway  -->  Lambda nexo-orders-api  -->  DynamoDB nexo-orders
    POST /orders  -->  API Gateway  -->  Lambda nexo-orders-api  -->  DynamoDB nexo-orders

  Admin (/admin tab Pedidos)
    GET /admin/orders          -->  API Gateway  -->  Lambda  -->  DynamoDB (Scan)
    PUT /admin/orders/{id}     -->  API Gateway  -->  Lambda  -->  DynamoDB (Update)

  Todos los requests llevan: Authorization: <JWT token de Cognito>
  API Gateway valida el token con el Cognito Authorizer antes de invocar el Lambda.
\end{lstlisting}

\subsection{Funciones --- src/lib/orders.ts}

\begin{table}[H]
\centering
\begin{tabular}{@{}ll@{}}
\toprule
\textbf{Función} & \textbf{Descripción} \\
\midrule
\texttt{getMyOrders()} & GET \texttt{/orders} --- pedidos del usuario autenticado \\
\texttt{addOrder(data)} & POST \texttt{/orders} --- crea un nuevo pedido \\
\texttt{getAllOrders()} & GET \texttt{/admin/orders} --- todos los pedidos (solo admin) \\
\texttt{updateOrderStatus(orderId, status)} & Atajo para actualizar solo el estado \\
\texttt{updateOrder(orderId, data)} & PUT \texttt{/admin/orders/\{id\}} --- actualiza status, peso o totalPagado \\
\bottomrule
\end{tabular}
\caption{API de pedidos --- src/lib/orders.ts}
\end{table}

\begin{nexobox}
Todas las llamadas a la API incluyen el token JWT en el header
\texttt{Authorization}, obtenido con \texttt{fetchAuthSession()} de Amplify.
El API Gateway valida el token con el Cognito Authorizer antes de invocar el Lambda.
\end{nexobox}

\subsection{Página del cliente --- /pedidos}

\begin{itemize}
  \item Lista de pedidos con badge de estado (En ruta / Atascado / Entregado)
  \item Muestra peso en libras y total pagado en USD si el admin los cargó
  \item Botón ``Agregar pedido'' abre formulario inline: tracking number (requerido) y descripción (opcional)
  \item Estado vacío con CTAs a \texttt{/cotizar} y WhatsApp
\end{itemize}

\subsection{Panel admin --- /admin tab Pedidos}

\begin{itemize}
  \item Tabla con todos los pedidos de todos los clientes
  \item Columnas: Cliente, Email, Tracking, Descripción, Fecha, Peso (lb), Total (\$), Estado
  \item Estado editable con \texttt{<select>} inline, guarda al cambiar
  \item Peso y Total: \texttt{<input type="number">} que guardan al perder foco (\textit{auto-save on blur})
\end{itemize}

% ==========================================================
\section{Panel de Administración}

\subsection{Protección por grupo Cognito}

\begin{lstlisting}[language=bash, caption={Verificacion de rol admin en /admin/page.tsx}]
const session = await fetchAuthSession()
const groups = session.tokens?.idToken
  ?.payload['cognito:groups'] as string[] ?? []
if (!groups.includes('admin')) {
  router.replace('/casillero')
  return
}
\end{lstlisting}

Los usuarios que no pertenecen al grupo \texttt{admin} de Cognito son redirigidos
a \texttt{/casillero} aunque conozcan la URL del panel.

\subsection{Tabs del dashboard}

\begin{table}[H]
\centering
\begin{tabular}{@{}lll@{}}
\toprule
\textbf{Tab} & \textbf{Datos} & \textbf{Fuente} \\
\midrule
Usuarios & Lista de cuentas de Cognito & \texttt{GET /admin/users} \\
Pedidos & Todos los pedidos del sistema & \texttt{GET /admin/orders} \\
\bottomrule
\end{tabular}
\caption{Tabs del panel de administración}
\end{table}

% ==========================================================
\section{Navbar y Menú de Usuario}

El \texttt{Navbar} detecta la sesión en cada cambio de ruta (\texttt{useEffect} sobre
\texttt{pathname}). Cuando hay sesión activa muestra el nombre del usuario con un
\textbf{dropdown} al hacer click:

\begin{itemize}
  \item \textbf{Mi casillero} → \texttt{/casillero}
  \item \textbf{Cuenta} → \texttt{/cuenta}
  \item \textbf{Cerrar sesión} → llama a \texttt{logoutUser()} y redirige a \texttt{/}
\end{itemize}

El dropdown se cierra automáticamente al hacer click fuera (listener
\texttt{mousedown} con \texttt{useRef}).

% ==========================================================
\section{Página de Cuenta --- /cuenta}

Permite al usuario autenticado gestionar su perfil:

\begin{itemize}
  \item \textbf{Datos editables:} nombre, apellido, móvil, teléfono
  \item \textbf{No editables desde aquí:} email ni contraseña
  \item \textbf{Guardar:} llama a \texttt{updateCurrentUser()} que usa
    \texttt{updateUserAttributes()} de Amplify
  \item \textbf{Eliminar cuenta:} doble confirmación. Llama a \texttt{deleteCurrentUser()}
    que ejecuta \texttt{deleteUser()} de Amplify, eliminando la cuenta de Cognito.
    Los pedidos en DynamoDB se \textbf{conservan} como registro interno.
\end{itemize}

% ==========================================================
\section{Variables de Entorno}

\subsection{Variables en Amplify Console}

\begin{table}[H]
\centering
\begin{tabular}{@{}ll@{}}
\toprule
\textbf{Variable} & \textbf{Valor} \\
\midrule
\texttt{NEXT\_PUBLIC\_COGNITO\_USER\_POOL\_ID} & \texttt{us-east-1\_Bphs6nkUf} \\
\texttt{NEXT\_PUBLIC\_COGNITO\_CLIENT\_ID} & \texttt{5fqqogr9akgphaabk9jk4tl3ap} \\
\texttt{NEXT\_PUBLIC\_ADMIN\_API\_URL} & \texttt{https://qduhjqlmfd.execute-api.us-east-1.amazonaws.com/prod} \\
\bottomrule
\end{tabular}
\caption{Variables de entorno en Amplify Console}
\end{table}

\subsection{Para desarrollo local}

Crear el archivo \texttt{.env.local} en la raíz del proyecto. Este archivo está en
\texttt{.gitignore} y \textbf{nunca debe subirse al repositorio}.

\begin{lstlisting}[language=bash, caption={.env.local --- solo para desarrollo}]
NEXT_PUBLIC_COGNITO_USER_POOL_ID=us-east-1_Bphs6nkUf
NEXT_PUBLIC_COGNITO_CLIENT_ID=5fqqogr9akgphaabk9jk4tl3ap
NEXT_PUBLIC_ADMIN_API_URL=https://qduhjqlmfd.execute-api.us-east-1.amazonaws.com/prod
\end{lstlisting}

% ==========================================================
\section{AWS Amplify Hosting}

\subsection{Configuración}

\begin{table}[H]
\centering
\begin{tabular}{@{}ll@{}}
\toprule
\textbf{Campo} & \textbf{Valor} \\
\midrule
App ID & \texttt{d3udh5pv9ejzs2} \\
Rama & \texttt{master} \\
Dominio & \texttt{nexocourier.com} \\
CloudFront URL & \texttt{dxykvodz7xo2w.cloudfront.net} \\
Región & \texttt{us-east-1} \\
SSL & AWS Certificate Manager (automático) \\
\bottomrule
\end{tabular}
\caption{Configuración de Amplify Hosting}
\end{table}

\subsection{Archivo amplify.yml}

\begin{lstlisting}[language=yaml, caption={amplify.yml}]
version: 1
frontend:
  phases:
    preBuild:
      commands:
        - npm ci --legacy-peer-deps
    build:
      commands:
        - npm run build
  artifacts:
    baseDirectory: .next
    files:
      - '**/*'
  cache:
    paths:
      - node_modules/**/*
      - .next/cache/**/*
\end{lstlisting}

% ==========================================================
\section{DNS y Dominio}

\begin{table}[H]
\centering
\begin{tabular}{@{}lllll@{}}
\toprule
\textbf{Tipo} & \textbf{Nombre} & \textbf{Valor} & \textbf{Proxy} & \textbf{TTL} \\
\midrule
CNAME & \texttt{www} & \texttt{dxykvodz7xo2w.cloudfront.net} & DNS only & Auto \\
CNAME & \texttt{@} & \texttt{dxykvodz7xo2w.cloudfront.net} & DNS only & Auto \\
\bottomrule
\end{tabular}
\caption{Registros DNS en Cloudflare}
\end{table}

\begin{nexobox}
\textbf{Importante:} Los registros CNAME deben estar en modo \textbf{DNS only}
(nube gris en Cloudflare), \textbf{no} en modo Proxy (nube naranja). Con el proxy
activo el certificado SSL de AWS no puede validarse.
\end{nexobox}

% ==========================================================
\section{Diseño y Sistema Visual}

\subsection{Paleta de colores}

\begin{table}[H]
\centering
\begin{tabular}{@{}lll@{}}
\toprule
\textbf{Token Tailwind} & \textbf{Hex} & \textbf{Uso} \\
\midrule
\texttt{space-black} & \texttt{\#0A0A0F} & Fondo principal de páginas \\
\texttt{midnight} & \texttt{\#0F0F18} & Fondo de cards y paneles \\
\texttt{cyan} & \texttt{\#00D4FF} & Color de marca, botones, CTAs \\
\texttt{ghost} & \texttt{\#E8E8F0} & Texto principal \\
\texttt{slate} & \texttt{\#6B7280} & Texto secundario, placeholders \\
\texttt{status-green} & \texttt{\#10B981} & Estados positivos \\
\texttt{status-yellow} & \texttt{\#F59E0B} & Advertencias, pendientes \\
\texttt{status-red} & \texttt{\#EF4444} & Errores y acciones destructivas \\
\bottomrule
\end{tabular}
\caption{Paleta de colores del sistema de diseño}
\end{table}

\subsection{Tipografía}

\begin{itemize}
  \item \textbf{Familia:} Inter (Google Fonts, cargada con \texttt{next/font/google})
  \item \textbf{Pesos:} 400, 500, 600, 700
  \item \textbf{Variable CSS:} \texttt{--font-inter}
\end{itemize}

% ==========================================================
\section{Comandos de Desarrollo}

\begin{lstlisting}[language=bash, caption={Comandos frecuentes}]
# Instalar dependencias (primera vez)
npm install --legacy-peer-deps

# Servidor de desarrollo local
npm run dev
# Disponible en: http://localhost:3000

# Verificar build antes de push
npm run build

# Deploy a produccion
git add <archivos>
git commit -m "tipo: descripcion"
git push
# Amplify detecta el push y despliega en ~3 minutos
\end{lstlisting}

% ==========================================================
\section{Repositorio y Organización AWS}

\begin{table}[H]
\centering
\begin{tabular}{@{}ll@{}}
\toprule
\textbf{Campo} & \textbf{Valor} \\
\midrule
Repositorio & \texttt{https://github.com/nexo-git/nexocode} \\
Rama principal & \texttt{master} \\
Organización GitHub & \texttt{nexo-git} \\
Usuario de commits & \texttt{AleKai-bot} \\
\midrule
Cuenta AWS prod & \texttt{nexo-prod} (769953010359) \\
Región de todos los recursos & \texttt{us-east-1} \\
\bottomrule
\end{tabular}
\caption{Repositorio y cuentas AWS}
\end{table}

\section*{Glosario}

\begin{description}
  \item[\textbf{App Router}] Sistema de enrutamiento de Next.js 14 basado en la
    estructura de carpetas dentro de \texttt{src/app/}.
  \item[\textbf{Auto-save on blur}] Patrón de UX donde un campo guarda su valor
    automáticamente cuando el usuario hace click fuera de él, sin botón de guardar.
  \item[\textbf{Double opt-in}] Flujo de registro que requiere confirmar la dirección
    de correo con un código antes de activar la cuenta.
  \item[\textbf{JWT}] JSON Web Token. Formato de token firmado que Cognito emite
    tras la autenticación exitosa.
  \item[\textbf{localStorage}] Almacenamiento del navegador donde Amplify v6 guarda
    los tokens de sesión. A diferencia de las cookies, no es accesible desde el
    Edge Runtime de Next.js.
  \item[\textbf{NEXT\_PUBLIC\_}] Prefijo de Next.js que expone variables de entorno
    al navegador del cliente.
  \item[\textbf{Cognito Authorizer}] Componente de API Gateway que valida el JWT
    de cada request antes de invocar el Lambda, extrayendo los claims automáticamente.
  \item[\textbf{Voseo}] Forma pronominal del español costarricense. Parte de la
    voz de marca nexo en toda la interfaz.
\end{description}

\vfill
\begin{center}
  \begin{tikzpicture}
    \fill[nexoBlack] (0,0) rectangle (\textwidth, 0.8cm);
    \node[anchor=west, text=nexoSlate, font=\small]
      at (0.3cm, 0.4cm) {nexo courier --- Confidencial --- Uso Interno --- v2.0};
    \node[anchor=east, text=nexoCyan, font=\small\itshape]
      at (\textwidth-0.3cm, 0.4cm) {Tu puente entre dos mundos.};
  \end{tikzpicture}
\end{center}

\end{document}
